\section{The retained inverse fibre as an explicit determinantal family}
\label{sec:fc-determinantal}

All rings in the algebraic statements are over a specified field $k$ of
characteristic zero. The bit statements use $k=\mathbb Q$ and signed binary
integers. No operation changes the coefficients of the original map.
The third source coordinate is named $w$; the renaming from the source's
$t$ is the identity $(x,y,t)\mapsto(x,y,w)$. The source expressions were
read in full from the parent satellite and the independent fluid extension;
their paths and SHA--256 hashes are in \texttt{source\_hashes.json}.
Every calculation used below is reproved here and replayed by the included
exact checker.

\subsection{The original map and its localized coordinate-ring isomorphism}

Put $A=k[a,b,c]$, $B=k[x,y,w]$, and define $F^*:A\to B$ by
\begin{align*}
 F^*(a)=F_1&=(1+xy)^3w+y^2(1+xy)(4+3xy),\\
 F^*(b)=F_2&=y+3x(1+xy)^2w+3xy^2(4+3xy),\\
 F^*(c)=F_3&=2x-3x^2y-x^3w.
\end{align*}
The determinant retains the source orientation $(x,y,w)$ and target
orientation $(a,b,c)$ and equals $-2$. Here is a direct proof. On $x\ne0$
set $v=1+xy$, $h=2-3xy-x^2w$, $A_0=v^2+v-v^3h$ and
$B_0=4v+2-3v^2h$. Their inverse coordinate formulas are
$y=(v-1)/x$ and $w=(5-3v-h)/x^2$, so these are actual coordinates
on the stated open set. Substitution in the original expressions yields
$F=(A_0/x^2,B_0/x,xh)$ and
\[
 \det\frac{\partial(x,v,h)}{\partial(x,y,w)}=-x^3,
 \qquad
 \det\frac{\partial F}{\partial(x,v,h)}
 =x^{-3}\det\begin{pmatrix}
 -2A_0&2v+1-3v^2h&-v^3\\
 -B_0&4-6vh&-3v^2\\h&0&1
 \end{pmatrix}=2x^{-3}.
\]
Expansion of the last numerator gives
\[
 (4-6vh)(-2v^2-2v+3v^3h)
 +(2v+1-3v^2h)(4v+2-6v^2h)=2.
\]
Multiplication proves $\det DF=-2$ on $x\ne0$. Both sides are
polynomials, and their equality in the localization of the domain $B$
proves equality in $B$ as well.

Retain the exact polynomial and derivative
\[
 P=cr^3-2r^2+br-2a,\qquad P'=3cr^2-4r+b,\qquad \alpha=P'/2.
\]
Define the $A$-algebra
\[
 E=\bigl(A[r]/(P)\bigr)[1/P'].
\]
There is an isomorphism of $A$-algebras $\Phi:E\to B[1/x]$ given by
\[
 (a,b,c,r,1/P')\longmapsto
 (F_1,F_2,F_3,y+1/x,x/2).
\]
Its inverse $\Psi:B[1/x]\to E$ is specified on every generator by
\[
 x\longmapsto 2/P',\quad
 y\longmapsto r-P'/2,\quad
 w\longmapsto 5(P'/2)^2-3r(P'/2)-c(P'/2)^3,
 \quad x^{-1}\longmapsto P'/2.
\]
To prove these assertions before quotienting, introduce independent
coordinates $(r,\alpha,c)$ with $\alpha\ne0$ and put
\[
 H(r,\alpha,c)=
 \left(\alpha^{-1},r-\alpha,5\alpha^2-3r\alpha-c\alpha^3\right).
\]
Solving $F_3=c$ after setting $r=y+1/x$, $\alpha=1/x$ gives this
formula and its inverse exactly. Direct substitution gives
\[
 F\circ H=(r^2+r\alpha-cr^3,\ 4r+2\alpha-3cr^2,\ c).
\]
The equation $F_2=b$ says $2\alpha=P'$, and then
$2(F_1-a)=P$. Conversely those two relations give $F_1=a,F_2=b$.
They show that the displayed maps respect $P=0$, the inverse elements,
and the $A$-algebra structure. The recovery equations
$r=y+1/x$, $\alpha=1/x$ show both compositions are the identity on
all their generators. This proves the ring isomorphism, including its
excluded divisor $P'=0$.

The morphism $F^*$ is injective. Indeed $A[r]/(P)$ is isomorphic to
$k[b,c,r]$ by eliminating $a=(cr^3-2r^2+br)/2$; both substitution maps
are inverse on the generators. A nonzero
$h(a,b,c)=\sum_{j=0}^d h_j(b,c)a^j$, $h_d\ne0$, has, after this
substitution, highest $r$ coefficient $h_d(b,c)c^d/2^d$ at degree $3d$,
which is nonzero in $k[b,c]$. Thus $A$ injects into the quotient and,
since $P'\ne0$ there, into $E$ and then $B[1/x]$. Consequently it
also injects into $B$. This proves dominance without treating it as an
assumption.

\subsection{Every specialized fibre and the finite free locus}

For $(a_0,b_0,c_0)\in k^3$, set
\[
 B_{a_0,b_0,c_0}=k[x,y,w]/(F_1-a_0,F_2-b_0,F_3-c_0),
 \qquad E_{a_0,b_0,c_0}=k[r,1/P']/(P),
\]
where $P$ and $P'$ here have the specialized coefficients.
If $c_0\ne0$, then $xh=c_0$ in the fibre ring, so $x$ is a unit
with inverse $h/c_0$. The preceding isomorphism therefore proves
$B_{a_0,b_0,c_0}\simeq E_{a_0,b_0,c_0}$ on the entire fibre.
If $c_0=0$, the entire fibre ring is instead
\[
 B_{a_0,b_0,0}\simeq k\ \times\ E_{a_0,b_0,0}.
\]
Here is a proof that includes its scheme structure. Work first in
$R=k[x,y,w]/(F_1-a_0,F_2-b_0)$. The polynomials $x,h$ generate the
unit ideal since
$1=h/2+x(3y+xw)/2$. Their ideals therefore have intersection $(xh)$:
if $z=xu=hv$, multiplication by the displayed expression for $1$
shows $z\in(xh)$. The map
$R/(xh)\to R/(x)\times R/(h)$ is injective by that intersection;
given classes modulo the two ideals, their lifts combine using the
same two summands of $1$ to prove surjectivity. This proves the asserted
product decomposition. The first factor is $k$ because
$F(0,y,w)=(w+4y^2,y,0)$, with point $(0,b_0,a_0-4b_0^2)$.
In the second factor, $x(3y+xw)=2$ makes $x$ a unit and the chart
isomorphism identifies it with $E_{a_0,b_0,0}$.

Over an algebraic closure, write $P=d\prod_j(r-\rho_j)^{m_j}$,
with distinct $\rho_j$. The coprime factors yield a product of the
rings $\bar k[r]/((r-\rho_j)^{m_j})$: the elementary two-ideal
argument just given applies repeatedly. When $m_j=1$, $P'$ has
nonzero value at $\rho_j$ and is a unit in that factor. When $m_j>1$,
$P'$ belongs to its nilpotent ideal $(r-\rho_j)$, so inverting $P'$
turns that factor into the zero ring. Thus $E$ retains exactly the
simple roots, each with its reduced point. The complete finite inverse
points are
\[
 \left(2/P'(\rho_j),\rho_j-P'(\rho_j)/2,
 5(P'(\rho_j)/2)^2-3\rho_jP'(\rho_j)/2-c_0(P'(\rho_j)/2)^3\right)
\]
for $m_j=1$, together with the boundary point when $c_0=0$.
The original polynomial, its multiplicities, and its excluded roots
remain recorded by $(P,P')$ even when the localization removes a factor.

The coefficient discriminant is exactly
\[
 D=4(b^2-cb^3+18abc-16a-27a^2c^2).
\]
For direct verification, expansion of the determinant
\[
 \det\begin{pmatrix}
 c&-2&b&-2a&0\\0&c&-2&b&-2a\\
 3c&-4&b&0&0\\0&3c&-4&b&0\\0&0&3c&-4&b
 \end{pmatrix}=-cD
\]
gives the stated polynomial. This determinant is the resultant because
its matrix represents $(u,v)\mapsto uP+vP'$ on the bases of
polynomials of degrees less than $2$ and $3$, respectively, with the
indicated descending coefficient order. Over a field it vanishes
exactly when $P,P'$ share a nonconstant factor: a nonzero kernel gives
$uP=-vP'$ and the coprime case forces $P\mid v$, impossible at
$\deg v<3$; a common factor gives a kernel vector by dividing both
polynomials by it. This also proves the simple-root interpretation.

Let $A_U=A[1/(cD)]$. Then $A_U[r]/(P)$ is free with basis
$(1,r,r^2)$. Polynomial division by the unit leading coefficient $c$
gives a spanning remainder of degree at most two, and a nonzero multiple
of $P$ has degree at least three, proving independence. No replacement
of $P$ is made here. The original relation gives the multiplication matrix
\[
 C=\begin{pmatrix}
 0&0&2a/c\\1&0&-b/c\\0&1&2/c
 \end{pmatrix},\qquad
 N=3cC^2-4C+bI_3.
\]
Its columns are exactly the coefficients of multiplying the ordered
basis by $r$ and $P'$, respectively. Direct determinant expansion gives
$\det N=-D/c$, so
$N^{-1}=-(c/D)\operatorname{adj}(N)$. This explicit inverse proves that
$P'$ is a unit in the quotient, and hence
\[
 B\otimes_A A_U\simeq A_U[r]/(P).
\]
This is a finite free rank-three morphism. Its unramified behavior can
also be proved without omitting the derivative: for a square-zero ideal
$I$ of an $A_U$-algebra $T$, a root $r_0$ in $T/I$ has a lift
$\widetilde r$ in $T$, and the unique root lift is
$\widetilde r-P(\widetilde r)/P'(\widetilde r)$.
The denominator is a unit because its reduction is; Taylor expansion
has no terms of degree two in $I$. If two lifts differ by $e\in I$,
subtraction gives $P'(\widetilde r)e=0$, so $e=0$. This is the
unique infinitesimal lifting property, establishing the finite etale
claim on precisely $U$.

At $c_0=0$, the polynomial still has degree two and its quadratic
discriminant is $b_0^2-16a_0$, while $D=4(b_0^2-16a_0)$. The
boundary point remains necessary. In particular the fibre at
$(-1/4,0,0)$ consists exactly of
\[
 (1,-3/2,13/2),\quad(-1,3/2,13/2),\quad(0,0,-1/4).
\]
At $c_0\ne0$ the cubic is a triple-root polynomial exactly at
$b_0=4/(3c_0)$, $a_0=4/(27c_0^2)$, with root $2/(3c_0)$.
Indeed comparing $c_0(r-\rho)^3$ with all four coefficients gives
these three equations and conversely verifies them. Its entire finite
fibre is empty since its only root has $P'=0$ and $c_0\ne0$ excludes
$x=0$.

\subsection{Generic degree, the absence of a rational inverse, and nonfiniteness}

In $k[a,b,c,r]$, $(P)$ is a prime ideal since its quotient was explicitly
identified with the domain $k[b,c,r]$. Thus $P$ is irreducible there.
It is primitive as a polynomial over $k[a,b,c]$ in $r$, since its
coefficient $-2$ is a unit. The primitive-polynomial argument proves
irreducibility over $K=k(a,b,c)$: clearing denominators in a proposed
factorization and cancelling the common prime factors of the coefficients
would produce a factorization in the unique-factorization ring
$k[a,b,c][r]$. To justify that cancellation, the product of two primitive
polynomials is primitive, since reduction modulo any prime leaves two
nonzero polynomials over a domain and hence a nonzero product. This
proves the required primitive-polynomial argument.

The chart identities now give
\[
 k(x,y,w)=k(F_1,F_2,F_3)(y+1/x),\qquad
 [k(x,y,w):k(F_1,F_2,F_3)]=3.
\]
Both containments in the first equality follow from the displayed
inverse formulas; irreducibility of the degree-three polynomial proves
the degree. A rational inverse would make the two fields equal, so none
exists. There is also no rational section over a dense target open set:
if its $x$ coordinate were identically zero then its $F_3$ coordinate
would be zero, contrary to the generic variable $c$; otherwise its
$r=y+1/x$ would be an element of $K$ satisfying the irreducible cubic.
These proofs rule out those exact rational maps and retain the
rank-three algebraic correspondence constructed above.

The morphism $F$ is not finite. Define the following homomorphism to
$k[t,t^{-1}]$ by its values on the original coordinates:
\[
 x=t^{-1},\quad y=-3t/2,\quad w=13t^2/2.
\]
Direct substitution gives $(F_1,F_2,F_3)=(-t^2/4,0,0)$.
If $B$ were finite over $A$, then multiplication by $x$ on finitely
many $A$-module generators would give a monic relation for $x$:
write $xm_i=\sum_j a_{ij}m_j$ and multiply the matrix $xI-(a_{ij})$
by its adjugate; its monic determinant annihilates every generator and
therefore $1\in B$. Under the displayed homomorphism such a relation
would become
\[
 t^{-m}+\sum_{j=0}^{m-1}q_j(-t^2/4,0,0)t^{-j}=0.
\]
Multiplication by $t^m$ gives a polynomial with constant term $1$,
while every summand in the sum is divisible by $t$. This is impossible.
Thus nonfiniteness is proved directly from the original map and is
compatible with its finite free restriction to $U$.

\subsection{The projective binary cubic restores the entire boundary sheet}

There is a stronger global presentation retaining the $x=0$ sheet as a
projective root. Define
\[
 Z=\{cR^3-2R^2T+bRT^2-2aT^3=0\}\subset\mathbb P^1_A,
 \qquad [R:T]=[1+xy:x].
\]
The last expression defines a morphism on the entire source because
$1=(1+xy)-yx$ proves its two homogeneous coordinates generate the
unit ideal. Substitution of $F_1,F_2,F_3$ in the homogeneous equation
gives zero: on $x\ne0$ it is $x^3P(y+1/x)=0$, and equality in the
localization of the source polynomial domain proves the polynomial
identity everywhere.

Let $Z^{\rm s}$ be the open simple-root locus: on $T\ne0$ its
coordinates are $r=R/T$ and it is defined by $P=0$, $P'\ne0$; on
$R\ne0$ use the distinct root coordinate $h=T/R$ and put
\[
 Q(h)=c-2h+bh^2-2ah^3,\qquad d=Q'(h)=-2+2bh-6ah^2.
\]
Its simple-root condition is $d\ne0$. On the overlap $h=1/r$,
$Q(h)=h^3P(1/h)$ and, on that zero scheme, $d=-hP'(1/h)$.
Since $h$ is a unit on the overlap, the two simple-root conditions
agree exactly.

The morphism from the source to $Z$ is an isomorphism onto $Z^{\rm s}$.
On the first chart this is the already proved inverse-root isomorphism.
Here is the entire second chart, including $x=0$:
\[
 x=-2h/d,\qquad y=b-3ah,\qquad
 w=-ad^3/8-y^2d^2/4+3y^2d/2.
\]
To prove it, first keep $a,b,h$ independent and define $d,x,y,w$ by
these displayed formulas with $d\ne0$. Substitution and expansion give
\[
 1+xy=-2/d,\qquad
 F_1=a,\qquad F_2=b,\qquad F_3=2h-bh^2+2ah^3=c-Q(h).
\]
For example $1+xy=(d-2hb+6ah^2)/d=-2/d$ directly. For the other
formulas one may insert that identity into the original $F_i$:
with $v=-2/d$, the first is
$v^3w+y^2v(1+3v)=a$, since $4+3xy=1+3v$;
the second is
$y+3xv^2w+3xy^2(1+3v)=y+3ha=b$.
The third follows by $F_3=x(5-3v-x^2w)$, inserting the displayed
$w$ and $y=b-3ah$, and expanding to $2h-bh^2+2ah^3$.
Thus the formulas define the inverse map on $Q=0$, $d\ne0$.
Conversely on the source chart $v=1+xy\ne0$, set $h=x/v$,
$a=F_1$, $b=F_2$, $c=F_3$. The polynomial identity gives $Q=0$;
the two original equations for $F_1,F_2$ give
$y=b-3ah$ by subtracting $3hF_1$ from $F_2$.
Consequently
$1+xy=1+x(b-3ah)=1+hv(b-3ah)$ gives
$v(1-bh+3ah^2)=1$, equivalently $d=-2/v$.
Then $x=hv=-2h/d$, and solving the original $F_1=a$ for $w$
gives exactly the stated $w$. These recover every generator, proving
both compositions. The source charts $x\ne0$ and $v\ne0$ cover it;
the two inverse maps agree on their overlap because both are inverse
to the same forward coordinates. This proves the global isomorphism.
At $h=0$, $Q=0$ requires $c=0$, $d=-2$, and the formulas give
$(x,y,w)=(0,b,a-4b^2)$. Thus the boundary point is retained as
$[R:T]=[1:0]$, with no discarded sheet.

The projective scheme $Z$ is finite locally free of rank three over
$A$. An elementary affine cover proves this. Evaluate its binary cubic
at the three constant points $[j:1]$, $j=0,1,-1$, obtaining
\[
 f_0=-2a,\quad f_1=c-2+b-2a,\quad f_{-1}=-c-2-b-2a.
\]
They generate the unit ideal since $f_1+f_{-1}-2f_0=-4$.
On the base open set $f_j\ne0$, the point $[j:1]$ is absent from $Z$.
Use the explicitly invertible homogeneous change
\[
 U=T,\quad V=R-jT,\qquad R=V+jU,\quad T=U.
\]
The removed point becomes $[U:V]=[1:0]$, so all of $Z$ on this
base open lies in the affine chart $V\ne0$ with $s=U/V$.
The affine coordinate formula is $R=1+js$, $T=s$; its polynomial is
\[
 c(1+js)^3-2(1+js)^2s+b(1+js)s^2-2as^3,
\]
whose leading coefficient is exactly $f_j$. It is a unit on this
base open, so division with the original polynomial gives the free
basis $1,s,s^2$ exactly as in the earlier cubic proof. The changes
are invertible projective linear maps and retain all original
coefficients. The three base opens cover, proving finite local
freeness and its rank.

The branch locus is exactly $D=0$. When $c\ne0$ this is the
resultant calculation already proved. When $c=0$, $[1:0]$ is a simple
root because $d(0)=-2$, and the remaining affine polynomial has degree
two with repeated root exactly when $b^2-16a=0$, equivalently $D=0$.
These arguments apply in the residue field at each base point and in
its algebraic closure. Consequently over $D\ne0$ every projective
root is simple, so $Z^{\rm s}=Z$ there and the original $F$ is finite
locally free of rank three on the larger open $\{D\ne0\}$, including
$c=0$. The unique root-lifting argument on the two charts proves it is
etale there. The algebra on $cD\ne0$ used previously gives its concrete
single-chart basis; that smaller open was sufficient and was not the
entire finite etale locus.

\subsection{A size-parameterized family with complete encoding}

For each integer $n\ge1$ retain independent variables
$(x_i,y_i,w_i)$ and targets $(a_i,b_i,c_i)$ for $1\le i\le n$.
Define $F^{[n]}$ to apply the preceding three original polynomials to
each block. This is a morphism $\mathbb A_k^{3n}\to\mathbb A_k^{3n}$
with coordinate homomorphism $a_i\mapsto F_1(x_i,y_i,w_i)$,
$b_i\mapsto F_2(x_i,y_i,w_i)$, $c_i\mapsto F_3(x_i,y_i,w_i)$.
Its Jacobian is block diagonal, so its determinant is exactly $(-2)^n$.
Its complete fibre algebra is the tensor product over $k$ of the
single-block fibre algebras just proved. This assertion follows by
presenting the tensor product as a polynomial ring on the disjoint sets
of generators modulo the union of their defining relations. It includes
every boundary factor at $c_i=0$.

Over $\{\prod_i D_i\ne0\}$ the block family is finite locally free
of rank $3^n$, by the preceding projective-cover proof and tensor products.
Over its subopen $U_n=\{\prod_i c_iD_i\ne0\}$, the algebra is finite free of
rank $3^n$ with the explicit basis
$\{\prod_i r_i^{e_i}:0\le e_i\le2\}$. Spanning follows by repeated
use of each original relation
$c_ir_i^3=2r_i^2-b_ir_i+2a_i$. Independence follows by taking tensor
products of the proven single-block bases: applying the corresponding
coordinate functionals to a putative relation extracts each coefficient.
The generic algebra is a domain since eliminating the $a_i$ identifies
its polynomial presentation with $k[b_i,c_i,r_i:1\le i\le n]$ before
localizing. A finite-dimensional domain over the target rational
function field is a field, because multiplication by a nonzero element
is an injective linear map and hence surjective. Thus the generic
function-field degree is also exactly $3^n$.

At the rational target with every block $(-1/4,0,0)$, the fibre contains
exactly the $3^n$ ordered choices of the three displayed rational points.
A format requiring one separate record for each point therefore has
output length at least $3^n$ symbols. In contrast, the tuple of equations
$(P_i,P_i')$, their localizations, and the specified boundary rule has
linear length in $n$ with fixed-size coefficients at this target. The
inequality follows from the mandated output format, and is not an
inference about the difficulty of obtaining that compact representation.

Here is an exact arithmetic-circuit model. A gate is binary $+,-$, or
$\times$, with fan-out allowed, and constants are $1,2,3,4$ for $F$.
Input nodes, constant nodes, and output wires are not counted as gates.
For one block the following list is a 26-gate circuit; each assignment
is exactly one gate and no listed coefficient is suppressed:
\[
\begin{array}{llll}
 g_1=xy,&g_2=1+g_1,&g_3=g_2g_2,&g_4=g_3g_2,\\
 g_5=3g_1,&g_6=4+g_5,&g_7=yy,&g_8=g_7g_6,\\
 g_9=g_2g_8,&g_{10}=g_4w,&g_{11}=g_{10}+g_9,&g_{12}=xg_3,\\
 g_{13}=g_{12}w,&g_{14}=3g_{13},&g_{15}=xg_8,&g_{16}=3g_{15},\\
 g_{17}=y+g_{14},&g_{18}=g_{17}+g_{16},&g_{19}=xx,&g_{20}=g_{19}x,\\
 g_{21}=g_{19}y,&g_{22}=3g_{21},&g_{23}=g_{20}w,&g_{24}=2x,\\
 g_{25}=g_{24}-g_{22},&g_{26}=g_{25}-g_{23}.&&
\end{array}
\]
The outputs are $(g_{11},g_{18},g_{26})$; substitution verifies all
three original formulas. Thus $F^{[n]}$ uses $26n$ gates. Its component
degrees remain $7,6,4$ per block. The tuple is uniformly constructible:
a loop on the binary index $i$ writes the same 26 instructions with
the input and gate indices of block $i$. Using binary indices, its
bit description has length $O(n\log(n+1))$, and the loop takes that
many elementary output operations.

For bit evaluation, encode each signed integer by a sign and its binary
magnitude, and assume magnitudes of the $3n$ inputs are less than $2^L$,
$L\ge1$. Every displayed intermediate has magnitude less than
$2^{7L+10}$. For example $|g_1|<2^{2L}$,
$|g_2|<2^{2L+1}$, $|g_6|<2^{2L+2}$,
$|g_{10}|<2^{7L+3}$, and $|g_9|<2^{6L+3}$;
the remaining products in the list have smaller exponent bounds, while
at most two sums require an extra bit. These inequalities also verify
the loose common bound for all gates. Grade-school multiplication on
$O(L)$ bits takes $O(L^2)$ bit operations by summing shifted partial
products; addition and subtraction take $O(L)$ operations by one pass
with carries. Evaluation therefore takes $O(n(L+1)^2)$ bit operations,
plus the circuit description or output indexing cost. This is a uniform
Boolean-time upper bound for this specified evaluation problem.

For target integers of magnitude less than $2^L$, retain the coefficient
tuples\[
 (c_i,-2,b_i,-2a_i),\qquad (3c_i,-4,b_i).
\]
Their entries have at most $L+3$ magnitude bits. Writing all polynomial
and localization equations
\[
 P_i(r_i)=0,\qquad z_iP_i'(r_i)-1=0
\]
therefore costs $O(n(L+1))$ bits and bit operations up to the index
encoding. The ring maps above identify $z_i$ with $1/P_i'$ and hence
$x_i=2z_i$ on its stated chart. Testing $c_i=0$ supplies the already
proved boundary rule. No root selection is performed by writing these
equations. When a rational root $r=p/q$, $q\ne0$, is supplied as a
certificate, with signed numerator and denominator lengths at most $M$,
integer cross multiplication verifies $P(r)=0$ and $P'(r)\ne0$ using
a fixed number of products of $O(L+M+1)$-bit integers. The coordinate
formulas then return exact rational numerators and denominators of
$O(L+M+1)$ bits in $O((L+M+1)^2)$ operations. Reduction to coprime
numerators and denominators is unnecessary: the output represents the
same rational numbers and retains the actual expression. This proves
the stated verification and lifting complexity while keeping the
selection of an algebraic embedding distinct from its full algebra.

\subsection{Exact determinant projections that retain every cubic coefficient}

A single retained polynomial has the following affine-linear matrix:
\[
 M(a,b,c,r)=\begin{pmatrix}
 r&-1&0&0\\0&r&-1&0\\0&0&r&-1\\-2a&b&-2&c
 \end{pmatrix},\qquad \det M=P.
\]
Expansion along the last row gives the four signed contributions
$-2a,br,-2r^2,cr^3$, respectively, proving the identity including signs.
Introduce one new variable $z$ and replace each constant entry by that
constant times $z$. This gives the homogeneous linear matrix
\[
 M^h=\begin{pmatrix}
 r&-z&0&0\\0&r&-z&0\\0&0&r&-z\\-2a&b&-2z&c
 \end{pmatrix},
 \quad \det M^h=cr^3-2z^2r^2+brz^2-2az^3=:P^h.
\]
Every monomial on the right has degree four; setting $z=1$ gives the
original polynomial exactly. On $z\ne0$ its relation to the affine
coordinates is the explicitly invertible substitution
$(a,b,c,r,z)\mapsto(a/z,b/z,c/z,r/z,z)$ with inverse
$(A,B,C,R,z)\mapsto(zA,zB,zC,zR,z)$, and
$P^h=z^4P(a/z,b/z,c/z,r/z)$. The divisor $z=0$ is newly added:
there $P^h=cr^3$. It is not identified with a finite original chart.

For $n$ blocks, the block-diagonal determinant $\prod_i P_i$ has
size $4n$. Its zero condition is a union of the individual hypersurfaces,
since a product in a field is zero exactly when a factor is zero.
It is not the simultaneous system $P_1=\cdots=P_n=0$ when $n\ge2$.
For example set all $r_i=b_i=c_i=0$, $a_1=0$, $a_2=1$, and for
$i>2$ set $a_i=1$. The product vanishes but $P_2=-2$.

The following determinant instead encodes every equation by a coefficient
map. Introduce independent selectors $t_1,\ldots,t_n$ and set
\[
 Q_n=\sum_{i=1}^n t_iP_i(r_i),\quad
 T_i=\begin{pmatrix}1&-r_i&0&0\\0&1&-r_i&0\\0&0&1&-r_i\\0&0&0&1\end{pmatrix},
 \quad v_i=\begin{pmatrix}-2a_i\\b_i\\-2\\c_i\end{pmatrix}.
\]
Let $T=\operatorname{diag}(T_1,\ldots,T_n)$,
$v=(v_1^{\mathsf T},\ldots,v_n^{\mathsf T})^{\mathsf T}$, and
$u=(t_1,0,0,0,\ldots,t_n,0,0,0)$. Then
\[
 \mathcal M_n=\begin{pmatrix}T&v\\-u&0\end{pmatrix},
 \qquad \det\mathcal M_n=Q_n,
 \qquad \operatorname{size}(\mathcal M_n)=4n+1.
\]
For a complete proof, write $T_i=I-N_i$. Since $N_i^4=0$,
$T_i^{-1}=I+N_i+N_i^2+N_i^3$, whose first row is
$(1,r_i,r_i^2,r_i^3)$. Thus $e_1^{\mathsf T}T_i^{-1}v_i=P_i$.
Left multiplication of $\mathcal M_n$ by the determinant-one matrix
$\left(\begin{smallmatrix}I&0\\uT^{-1}&1\end{smallmatrix}\right)$
gives $\left(\begin{smallmatrix}T&v\\0&uT^{-1}v\end{smallmatrix}\right)$.
As $\det T=1$, its determinant is $uT^{-1}v=Q_n$.
Every entry of $\mathcal M_n$ is affine-linear in the original variables
and selectors with coefficients in $\{0,1,-1,-2\}$.

For an arbitrary coefficient ring $R$, the map
$\iota:R^n\to R[t_1,\ldots,t_n]_1$,
$(f_1,\ldots,f_n)\mapsto\sum_i t_if_i$, is an $R$-module
isomorphism onto the homogeneous selector-degree-one part, with inverse
$Q\mapsto([t_1]Q,\ldots,[t_n]Q)$. Comparing the unique monomial
coefficients proves both compositions are identities. In particular
$(P_1,\ldots,P_n)=([t_1]\det\mathcal M_n,\ldots,[t_n]\det\mathcal M_n)$
as ideals in $k[a_i,b_i,c_i,r_i]$. For a specialized point over any
extension field, all fibre equations hold if and only if the determinant
is identically zero as a polynomial in the selectors. Specializing the
selectors first can lose information: when $t_1=t_2=1$, the choices
$P_1=2$, $P_2=-2$ give zero sum without either equation holding.
The coefficient map is an explicit reversible transport of the tuple,
not an asserted isomorphism between the fibre scheme and one
hypersurface in a larger space.

Replace each constant in $\mathcal M_n$ by that constant times $z$;
call the resulting homogeneous linear matrix $\mathcal M_n^h$.
Its determinant is
\[
 \det\mathcal M_n^h
 =z^{4n-4}\sum_{i=1}^n t_i
   (c_ir_i^3-2z^2r_i^2+b_ir_iz^2-2a_iz^3).
\]
To prove this without assigning an inverse at $z=0$, first localize at
$z$. The four diagonal entries of each $T_i$ become $z$, and the first
row of its inverse is $(z^{-1},r_iz^{-2},r_i^2z^{-3},r_i^3z^{-4})$.
The last vector becomes $(-2a_i,b_i,-2z,c_i)^{\mathsf T}$, and
$\det T=z^{4n}$. Applying the same determinant-one block multiplication
gives the displayed formula in $k[\text{variables},z,z^{-1}]$.
Both sides belong to the polynomial subring, whose localization is
injective, so the identity holds also at $z=0$. The factor
$z^{4n-4}$ is explicit: for $n>1$ the entire hyperplane $z=0$ is an
additional zero component of the padded determinant. At $z=1$ all
original fibre coefficients are recovered.

The determinantal projection is an exact coordinate-ring homomorphism.
With $d=4n+1$ let $k[X_{pq}:1\le p,q\le d]$ be the coordinate ring
of all $d\times d$ matrices. The homomorphism
$\vartheta_n:k[X_{pq}]\to k[a_i,b_i,c_i,r_i,t_i,z]$ sends
$X_{pq}$ to the displayed $(p,q)$ entry of $\mathcal M_n^h$.
It sends the universal determinant to the stated padded polynomial.
It is surjective: each $a_i,b_i,c_i,r_i,t_i$ is a nonzero scalar
multiple of a specified entry, and a diagonal entry gives $z$.
Its kernel is generated by the linear relations among those entries.
To prove the last claim constructively, for each target variable choose
one entry mapping to it after its indicated scalar division. Subtract
from every other $X_{pq}$ the same linear combination of these chosen
entries as its image. These linear differences lie in the kernel.
Modulo them the source is the polynomial ring on the chosen entries;
the induced map to the target polynomial ring has the inverse sending
each target variable to its chosen entry. Therefore the kernel is
exactly the ideal generated by those differences. This proves the map,
its image, its kernel and the degree $d$ of the determinant image.

Finally, Horner evaluation of the unaltered cubic uses seven gates:
$c_ir_i$, then addition of $-2$, multiplication by $r_i$, addition of
$b_i$, multiplication by $r_i$, separately $(-2)a_i$, and addition of
that last term. Computing all $t_iP_i$ and their sum uses $9n-1$ gates.
Thus $Q_n$ has an explicitly uniform linear-size arithmetic circuit and
a linear-size affine determinantal representation. The homogeneous padded
form has a linear-size determinant representation as well. These are
constructive upper bounds for a family derived from the original fibres.
They retain their generic degree $3^n$ and their original exceptional
loci. They prove no reduction from arbitrary Boolean satisfiability to
these independent blocks, and they assert no Boolean complexity equality
or separation. The concrete complexity transport established here is
the coefficient-recoverable determinant projection and the fully specified
uniform bit algorithms above.
